\subsection{Linear Functions} 
Linear equations, expressions, and functions are some of the easiest ones to work with. We can define them in various ways.
\begin{definition}
A \term{linear function} is one that can be written in the form $$f(x)=mx+b$$
\end{definition}
The main issue with this definition is with the phrase ``\makeblank[1.75in].'' When you see a function, how can you tell whether or not it's linear?
\begin{definition}
A \term{linear expression} is one in which the only operations are:
\begin{itemize}
\item addition and subtraction
\item multiplication and division \emph{only by numbers}
\begin{itemize}
\item Multiplying variables together and dividing by variables is not allowed.
\end{itemize}
\end{itemize}
A \term{linear equation} is one in which the expressions on each side are linear. A \term{linear function} is one that takes the form $f(x)=\mbox{[linear expression]}$.
\end{definition}
\begin{Example}
Each expression below is not linear. Explain why.
\begin{itemize}
\item $x+\frac{1}{x}$\makeblanknewf
\item $\sqrt{x-5}$\makeblanknewf
\item $x+2xy+y$\makeblanknewf
\item $(x-3)(x+5)$\makeblanknewf
\end{itemize}
\end{Example}
\begin{Note}
Depending on the context, we may or may not think of an expression with no variables (a \makeblank[1in] expression) as linear. For our purposes, we will say that linear expressions must contain one or more variables.
\end{Note}
\begin{Note}
Technically, the functions we call \term{linear} at this level are more properly referred to as \term{affine}. We will ignore this issue in this course outside of this note.
\end{Note}